\documentclass[11pt,a4paper]{article}
\usepackage[utf8]{inputenc}
\usepackage[T1]{fontenc}
\usepackage[english]{babel}
\usepackage{amsmath,amssymb,amsthm}
\usepackage{geometry}
\geometry{margin=2.5cm}
\usepackage{hyperref}
\usepackage{booktabs}
\usepackage{xcolor}
\usepackage{mdframed}
\usepackage[expansion=false]{microtype}
\usepackage{enumitem}
\usepackage{array}
\usepackage{graphicx}
\usepackage{caption}
\captionsetup{font=small,labelfont=bf,format=plain}
\hypersetup{colorlinks=true, linkcolor=blue!50!black, urlcolor=blue!50!black}

% Epistemic boxes
\newmdenv[backgroundcolor=green!10, linecolor=green!50!black, linewidth=1pt,
  innertopmargin=6pt, innerbottommargin=6pt, innerleftmargin=8pt, innerrightmargin=8pt,
  skipabove=8pt, skipbelow=8pt]{rigorousbox}
\newmdenv[backgroundcolor=yellow!15, linecolor=orange!70!black, linewidth=1pt,
  innertopmargin=6pt, innerbottommargin=6pt, innerleftmargin=8pt, innerrightmargin=8pt,
  skipabove=8pt, skipbelow=8pt]{heuristicbox}
\newmdenv[backgroundcolor=red!8, linecolor=red!50!black, linewidth=1pt,
  innertopmargin=6pt, innerbottommargin=6pt, innerleftmargin=8pt, innerrightmargin=8pt,
  skipabove=8pt, skipbelow=8pt]{negativebox}
\newmdenv[backgroundcolor=blue!8, linecolor=blue!60!black, linewidth=1pt,
  innertopmargin=6pt, innerbottommargin=6pt, innerleftmargin=8pt, innerrightmargin=8pt,
  skipabove=8pt, skipbelow=8pt]{openproblem}
\newmdenv[backgroundcolor=violet!8, linecolor=violet!60!black, linewidth=1pt,
  innertopmargin=6pt, innerbottommargin=6pt, innerleftmargin=8pt, innerrightmargin=8pt,
  skipabove=8pt, skipbelow=8pt]{discoverybox}

\theoremstyle{plain}
\newtheorem{theorem}{Theorem}[section]
\newtheorem{proposition}[theorem]{Proposition}
\newtheorem{corollary}[theorem]{Corollary}
\newtheorem{conjecture}[theorem]{Conjecture}
\newtheorem{lemma}[theorem]{Lemma}

\theoremstyle{definition}
\newtheorem{definition}[theorem]{Definition}
\newtheorem{remark}[theorem]{Remark}
\newtheorem{hypothesis}[theorem]{Hypothesis}
\newtheorem{observation}[theorem]{Observation}

\title{\large $p$-adic Distribution of the Minimal Offset $q_{\min}$\\
in Parametric Prime Families:\\[3pt]
\Large Bateman--Horn Predictions, Conditional Theorems,\\
and an Anti-Persistence Sign-Reversal Phenomenon}
\author{Hassane Bakkaoui\\\small Independent Researcher\\\small bakkahassa@hotmail.com}
\date{Working draft v0.1 --- May 2026\\\small Sequel to Papers I and II (April 2026)}

\newcommand{\qmin}{q_{\min}}
\newcommand{\Z}{\mathbb{Z}}
\newcommand{\N}{\mathbb{N}}
\newcommand{\Q}{\mathbb{Q}}
\newcommand{\R}{\mathbb{R}}
\newcommand{\Prim}{\mathbb{P}}
\newcommand{\E}{\mathbb{E}}
\newcommand{\Prb}{\mathbb{P}}
\newcommand{\PiP}{\pi_{\mathrm{perm}}}
\newcommand{\PiF}{\pi_{\mathrm{forb}}}
\newcommand{\eps}{\varepsilon}
\newcommand{\Eclass}{\mathcal{E}}

\begin{document}
\maketitle

\begin{abstract}
We study the distribution modulo small primes $\ell$ of the minimal offset
$\qmin(k, m, \eps)$ for the parametric prime family
$p = km(m+1) + \eps + 2kq$, with $k \ge 1$ and $\eps$ admissible
(\textit{i.e.}, $\gcd(\eps, 2k) = 1$). This completes a trilogy of papers on
this family: Paper~I (April 2026) characterised the marginal law
$\mathbb{E}[|\qmin|] \sim \log m / C_k$; Paper~II (April 2026, revised)
established the conditional law $\mathbb{E}[|q| \mid m] = m/4$ and excluded a
spurious spectral connection to $\zeta(s)$; the present paper analyses the
\emph{joint modular} layer, completing the description with a third axis.

\smallskip
\noindent Our principal contributions are:

\smallskip
(i) An \textbf{explicit forbidden residue} $q^*(k, m, \eps; \ell) \in \Z/\ell\Z$
(Proposition~\ref{prop:forbidden}), giving rise to the \textbf{permission frequency}
$\PiP(r; k, \ell)$.

\smallskip
(ii) A \textbf{conditional theorem} (Theorem~\ref{thm:main}): under H1$'$ and
H3$'$, $F_N(r; k, \ell) = \PiP(r; k, \ell)/(\ell-1) + O(N^{-1/2})$. Validated
empirically at $N = 10^6$ for $k \in \{3, 7, 15, 21\}$: mean Pearson correlation
prediction/observation $r = 0.99$ for small $\ell$, with effective constant
$\bar C \approx 0.5$ (5$\times$ smaller than the triangle-inequality bound).

\smallskip
(iii) An \textbf{unconditional theorem} (Theorem~\ref{thm:degenerate}) for
$\ell \mid 2k$: $F_N(r) = 1/\ell + O(N^{-1/2})$ unconditionally; verified across
all 10 tested cases with $\chi^2 < 3.1$.

\smallskip
(iv) A \textbf{new phenomenon} (Section~\ref{sec:anti-persistence}): the lag-1
autocorrelation of $(q_n \bmod 2)$ follows a clean \textbf{linear law in $k$},
$\mathrm{ACF}_1 \approx 0.024\,k - 0.20$ ($R^2 = 0.97$ across 10 values of
$k \in [3, 21]$ at $N = 10^6$), with sign-reversal threshold at $k_0 \approx 8.4$.
The full spectral signature $\{\alpha(\ell)\}_{\ell=1}^{10}$ of the slopes per
lag is fitted by a \textbf{damped cosine} (period $T \approx 5.89$, damping
rate $\gamma \approx 0.39$) to within $R^2 = 0.999$. The phenomenon
\textbf{amplifies with the modulus} (factor $\sim 2\times$ between mod 2 and
mod 5 at large $k$). At $k=21$, mod 5: $\mathrm{ACF}_1 = +0.64$.

\smallskip
\textbf{Connection to Paper~II}: Hypothesis H3$'$ is a structural analogue of
Paper~II's H3 for the residue class. Theorem~\ref{thm:main} thereby provides
an \textbf{indirect quantitative test} of H3 via measurable distributions.
The empirical agreement is consistent with H3 holding in the regime tested,
but this constitutes \emph{empirical support}, not proof: the unconditional
barrier (parity problem; uniform PNT in arithmetic progressions of large
moduli) is unchanged.

\smallskip
\textbf{Status of results}: Theorem~\ref{thm:degenerate} is unconditional.
Theorem~\ref{thm:main}, Proposition~\ref{prop:h3-implies-h3prime}, and the
ACF conjectures are heuristic or conditional, with carefully tracked
hypotheses. We discuss the obstacles to unconditional proof in
Section~\ref{sec:connection}.

\smallskip
\textbf{MSC 2020}: 11N32, 11N13, 11Y11, 11M41, 62P99.\\
\textbf{Keywords}: prime numbers, parametric families, $p$-adic distribution,
Bateman--Horn conjecture, autocorrelation, anti-persistence.
\end{abstract}

\begin{heuristicbox}
\noindent\textbf{Epistemic key.} Throughout the paper, every result carries one of
five labels: \textbf{Rigorous} (unconditional, proved from classical theorems);
\textbf{Heuristic} (conditional on H1$'$ or H3$'$); \textbf{Discovery} (new
empirical finding requiring further theoretical work); \textbf{Negative}
(invalidation of a prior small-$N$ observation); \textbf{Open} (problem stated
but unsolved).
\end{heuristicbox}

\tableofcontents

\section{Introduction}
\label{sec:intro}

\subsection{From marginal to conditional to modular distributions}

The parametric prime family
\begin{equation}
p = km(m+1) + \eps + 2kq, \qquad m \in \N^*,\ \eps \in \Eclass_k,\ q \in \Z,
\label{eq:family}
\end{equation}
introduced in Paper~I~\cite{paper1} and refined in Paper~II~\cite{paper2}, captures
each prime $p > k+1$ as a triple $(m, \eps, q)$ with $|q|$ minimal. Papers~I and~II
analysed the \emph{size} of $q$:

\begin{itemize}[leftmargin=*]
\item Paper~I: $\E[|\qmin(k, m)|] \sim \log m / C_k$, validated at $N = 10^6$ for $k=21$
(heuristic, under H1+H2).
\item Paper~II: $\E[|q| \mid m] = m/4$ (Conditional Proposition~9.4), and the
universal constant $C_0(k) = 1/(4\sqrt{k})$ across $k \in \{3, 15, 21\}$.
\end{itemize}

The present paper analyses the \emph{joint modular structure} of $\qmin$: how is
$\qmin$ distributed modulo a small prime $\ell$? This is the natural next layer
in the description, complementing the marginal (Paper~I) and conditional
(Paper~II) layers with the modular layer.

\subsection{Pilot observation and main motivating question}

A pilot study at $N = 10^4$ for $k = 3$ revealed a striking dichotomy:
\begin{itemize}[leftmargin=*]
\item For $\ell \in \{2, 3\}$ (which divide $2k = 6$): the distribution of
$\qmin \bmod \ell$ is \emph{exactly uniform}.
\item For $\ell \in \{5, 7, 11, 13\}$: the distribution is \emph{strongly
non-uniform} ($p$-values $< 10^{-6}$), but the bias follows a precisely
predictable pattern.
\end{itemize}

The principal goal of this paper is to (a) characterise this pattern theoretically,
(b) prove (under reasonable hypotheses) a quantitative theorem, and (c) validate
it across $k \in \{3, 7, 15, 21\}$ at $N = 10^6$.

A second, unexpected discovery emerged from the validation campaign: a clean
\textbf{linear law} for the lag-1 autocorrelation of $(q \bmod 2)$ in terms of
$k$ (Figure~\ref{fig:acf-teaser}), with a sign-reversal threshold near
$k_0 \approx 8.4$. Tested across 10 values of $k \in [3, 21]$ at $N = 10^6$:
$\mathrm{ACF}_1 \approx 0.024\,k - 0.20$ ($R^2 = 0.97$).

\begin{figure}[h!]
\centering
\includegraphics[width=0.95\textwidth]{fig4_acf_reversal.png}
\caption{Discovery (preview): lag-1 autocorrelation of $(q_n \bmod 2)$ for
four representative values of $k$ at $N = 10^6$. \textbf{(A)} Sign reversal
between $k = 7$ (anti-persistence) and $k = 15$ (positive persistence).
\textbf{(B)} Empirical probability of same-parity transitions, reaching
$0.658$ at $k = 21$ (excess of $157{,}600$ over independence baseline). The
extended investigation across 10 values of $k$ reveals a linear law in $k$
(Figure~\ref{fig:acf-extended} and Section~\ref{sec:anti-persistence}).}
\label{fig:acf-teaser}
\end{figure}

\subsection{Main results}

\textbf{Theorem~\ref{thm:degenerate}} (rigorous). For $\ell \mid 2k$,
$F_N(r; k, \ell) = 1/\ell + O(N^{-1/2})$ unconditionally.

\medskip
\textbf{Theorem~\ref{thm:main}} (heuristic, under H1$'$+H3$'$). For
$\gcd(2k, \ell) = 1$, $F_N(r; k, \ell) = \PiP(r; k, \ell)/(\ell-1) + O(N^{-1/2})$.

\medskip
\textbf{Conjecture~\ref{conj:tight}}. The implicit constant in Theorem~\ref{thm:main}
is bounded by an absolute $C_{\mathrm{opt}} \approx 0.5$, independent of $k$.

\medskip
\textbf{Discovery~\ref{disc:acf}}. Lag-1 ACF sign reversal of $(q \bmod 2)$ between
$k = 7$ and $k = 15$.

\subsection{Organisation and reading guide}

The paper is organised into three logical layers:

\paragraph{Theoretical core (Sections~\ref{sec:framework} and~\ref{sec:connection}).}
Section~\ref{sec:framework} establishes the theoretical framework, including
Hypotheses H1$'$ and H3$'$ and the proofs of Theorems~\ref{thm:degenerate} and
\ref{thm:main}. Section~\ref{sec:connection} formally establishes the
implication H3 (Paper~II) $\Rightarrow$ H3$'$ via
Proposition~\ref{prop:h3-implies-h3prime}, with explicit error tracking, and
discusses at length the obstacles to unconditional proof.

\paragraph{Empirical validation (Sections~\ref{sec:validation-k3}
and~\ref{sec:universality}).}
Section~\ref{sec:validation-k3} validates Theorems~\ref{thm:main}
and~\ref{thm:degenerate} for $k = 3$ at $N = 10^6$, including the scaling
test of $1/\sqrt{N}$ across four decades and the verification of H3$'$ at the
stratum level. Section~\ref{sec:universality} extends the validation across
$k \in \{3, 7, 15, 21\}$, demonstrating universality of the theorems.

\paragraph{New phenomenon (Section~\ref{sec:anti-persistence}).}
A genuinely new empirical discovery---a damped-oscillation spectral signature
of the lag-$\ell$ autocorrelation of $(q_n \bmod 2)$ as a function of $k$---is
documented in detail across 10 values of $k$. The discovery is presented as a
phenomenon awaiting theoretical explanation (Open Problem~D).

\paragraph{Methodological and bibliographic matter
(Sections~\ref{sec:negative}--\ref{sec:methodology}).}
Section~\ref{sec:negative} addresses small-$N$ artefacts and methodology
caveats. Section~\ref{sec:openpb} lists six open problems
(\textsf{A}--\textsf{F}). Section~\ref{sec:methodology} provides the
computational methodology with cost estimates for reproducibility.

A summary epistemic table (Section~\ref{sec:epistemic}) classifies all results
by their status (rigorous, heuristic, conditional, conjectural).

\smallskip
\noindent\textbf{Reading suggestions.} A reader interested only in the rigorous
content may focus on Proposition~\ref{prop:forbidden} and
Theorem~\ref{thm:degenerate} (Section~\ref{sec:framework}). A reader interested
in the heuristic theory can read Section~\ref{sec:framework} in full and the
connection to Paper~II (Section~\ref{sec:connection}). A reader interested in
the empirical discovery alone can begin with
Section~\ref{sec:anti-persistence}, which is largely self-contained.

\section{Theoretical Framework}
\label{sec:framework}

\subsection{Notations and admissible classes}

Fix $k \ge 1$. The admissible classes for $\eps$ are
\[
\Eclass_k = \{\eps \in \Z : -k < \eps \le k,\ \gcd(\eps, 2k) = 1\}.
\]
By a standard application of Dirichlet's theorem~\cite{dirichlet}, every prime
$p > k+1$ admits a unique representation~\eqref{eq:family} with $|q|$ minimal
(possibly with ties resolved by convention; see Paper~II). Set
\[
\qmin(k, m, \eps) = \arg\min_{q \in \Z}\{|q| : km(m+1) + \eps + 2kq \in \Prim\}.
\]
We have $|\Eclass_k| = \varphi(2k)$. For the values studied,
$|\Eclass_3| = 2$, $|\Eclass_7| = 6$, $|\Eclass_{15}| = 8$, $|\Eclass_{21}| = 12$.

\subsection{The forbidden residue}

\begin{rigorousbox}
\begin{proposition}[Forbidden residue]\label{prop:forbidden}
Let $\ell$ be a prime with $\gcd(\ell, 2k) = 1$. For every triple $(k, m, \eps)$,
there exists a unique residue $q^*(k, m, \eps; \ell) \in \Z/\ell\Z$ such that
\[
q \equiv q^*(k, m, \eps; \ell) \pmod{\ell} \implies \ell \mid p_{k, m, \eps, q}.
\]
Explicitly:
\begin{equation}
q^*(k, m, \eps; \ell) \equiv -(2k)^{-1}\bigl(km(m+1) + \eps\bigr) \pmod{\ell}.
\label{eq:qstar}
\end{equation}
\end{proposition}
\end{rigorousbox}

\begin{proof}
The condition $p \equiv 0 \pmod \ell$ rewrites as
$2kq \equiv -(km(m+1) + \eps) \pmod \ell$. Since $\gcd(2k, \ell) = 1$, the inverse
$(2k)^{-1}$ exists and the equation has the unique solution~\eqref{eq:qstar}.
\end{proof}

\begin{corollary}[Permitted residues]\label{cor:permitted}
For each triple $(k, m, \eps)$, exactly $\ell - 1$ residues of $q$ modulo $\ell$
are \emph{permitted} (do not impose $\ell \mid p$).
\end{corollary}

\subsection{Aggregate forbid and permit frequencies}

\begin{definition}\label{def:perm}
The \emph{forbidden frequency} of residue $r \in \Z/\ell\Z$ is
\[
\PiF(r; k, \ell) = \frac{1}{\ell |\Eclass_k|}
\#\bigl\{(m_0, \eps) \in \Z/\ell\Z \times \Eclass_k : q^*(k, m_0, \eps; \ell) \equiv r\bigr\}.
\]
The \emph{permission frequency} is $\PiP(r; k, \ell) = 1 - \PiF(r; k, \ell)$.
\end{definition}

The identity $\sum_r \PiF(r) = 1$ (each pair $(m_0, \eps)$ forbids exactly one
residue) implies $\sum_r \PiP(r) = \ell - 1$.

\subsection{Hypotheses}

\begin{heuristicbox}
\begin{hypothesis}[H1$'$: equidistribution of $(m \bmod \ell, \eps)$]\label{hyp:H1prime}
Let $S_N = \{(m_n, \eps_n)\}_{n=1}^N$ be the admissible triples generating the
first $N$ primes of family~\eqref{eq:family} in increasing order of $p$. Define
$\nu_N(m_0, \eps) = \frac{1}{N} \#\{n \le N : m_n \equiv m_0 \pmod \ell, \eps_n = \eps\}$.
Then, uniformly in $(m_0, \eps)$,
\[
\nu_N(m_0, \eps) = \frac{1}{\ell |\Eclass_k|} + O(N^{-1/2}).
\]
\end{hypothesis}
\textbf{Status.} A consequence of Bombieri--Vinogradov at fixed modulus
$\ell |\Eclass_k|$. Empirically validated: maximum deviation $0.36\%$ at
$k=3$, $\ell=5$, $N=10^4$, all $z$-scores below $1.2\sigma$.
\end{heuristicbox}

\begin{heuristicbox}
\begin{hypothesis}[H3$'$: conditional uniformity on permitted residues]\label{hyp:H3prime}
For each pair $(m_0 \in \Z/\ell\Z, \eps \in \Eclass_k)$, conditionally on
$m_n \equiv m_0 \pmod \ell$ and $\eps_n = \eps$, the distribution of
$\qmin(k, m_n, \eps_n) \bmod \ell$ is uniform on the $\ell-1$ permitted residues:
\[
\Prb\!\left(\qmin \equiv r \pmod \ell \mid m \equiv m_0, \eps\right)
= \frac{\mathbb{1}[r \neq q^*(k, m_0, \eps; \ell)]}{\ell - 1}.
\]
\end{hypothesis}
\textbf{Status.} Heuristic; an analogue of Hypothesis H3 of Paper~II. Empirically
validated to $99.9999\%$ at $k=3$, $\ell=5$, $N = 10^6$ (1 violation in $10^6$
cases tested across 10 strata).
\end{heuristicbox}

\subsection{Main theorem (non-degenerate case)}

\begin{heuristicbox}
\begin{theorem}[$p$-adic distribution of $\qmin$]\label{thm:main}
Let $\ell$ be a prime with $\gcd(2k, \ell) = 1$. Under
Hypotheses~\ref{hyp:H1prime} and~\ref{hyp:H3prime}, for every $r \in \Z/\ell\Z$:
\begin{equation}
F_N(r; k, \ell) := \frac{1}{N} \#\{n \le N : \qmin(k, m_n, \eps_n) \equiv r \pmod \ell\}
= \frac{\PiP(r; k, \ell)}{\ell - 1} + O_{k, \ell}(N^{-1/2}).
\label{eq:main}
\end{equation}
\end{theorem}
\end{heuristicbox}

\begin{proof}
\emph{Step 1 (Decomposition).} By H3$'$ and total probability:
\[
F_N(r) = \sum_{m_0, \eps} \nu_N(m_0, \eps) \cdot
\frac{\mathbb{1}[r \neq q^*(k, m_0, \eps; \ell)]}{\ell - 1}.
\]

\emph{Step 2 (Target).} By Definition~\ref{def:perm},
\[
\frac{\PiP(r)}{\ell - 1}
= \frac{1}{\ell |\Eclass_k| (\ell - 1)} \sum_{m_0, \eps} \mathbb{1}[r \neq q^*].
\]

\emph{Step 3 (Difference).}
\[
F_N(r) - \frac{\PiP(r)}{\ell - 1}
= \frac{1}{\ell - 1} \sum_{m_0, \eps} \left[\nu_N(m_0, \eps) - \frac{1}{\ell |\Eclass_k|}\right]
\mathbb{1}[r \neq q^*].
\]

\emph{Step 4 (Triangle inequality).}
\[
\left|F_N(r) - \frac{\PiP(r)}{\ell - 1}\right|
\le \frac{1}{\ell - 1} \sum_{m_0, \eps} \left|\nu_N(m_0, \eps) - \frac{1}{\ell |\Eclass_k|}\right|.
\]

\emph{Step 5 (Conclusion).} By H1$'$, each summand is $O(N^{-1/2})$, with
$\ell |\Eclass_k|$ terms. Hence
$\left|F_N(r) - \PiP(r)/(\ell-1)\right| = O\bigl(\ell |\Eclass_k|/(\ell-1) \cdot N^{-1/2}\bigr)
= O_{k, \ell}(N^{-1/2})$.
\end{proof}

\begin{corollary}[Explicit constant]\label{cor:constant}
The implicit constant in Theorem~\ref{thm:main} is
\[
C(k, \ell) = \frac{\ell |\Eclass_k|}{\ell - 1} \cdot \sigma(k, \ell), \qquad
\sigma(k, \ell) = \sqrt{\frac{\ell |\Eclass_k| - 1}{\ell |\Eclass_k|}} \le 1.
\]
For $k = 3$, $\ell = 5$: $C(3, 5) \approx 2.37$.
\end{corollary}

\subsection{Degenerate case (rigorous)}

\begin{rigorousbox}
\begin{theorem}[Uniform distribution in the degenerate case]\label{thm:degenerate}
If $\ell \mid 2k$, then for every triple $(k, m, \eps)$ and every $q$:
$p_{k, m, \eps, q} \equiv \eps \pmod \ell$, which is non-zero modulo $\ell$ since
$\gcd(\eps, 2k) = 1$ and $\ell \mid 2k$. Therefore no residue of $q$ is forbidden,
and Theorem~\ref{thm:main} degenerates to:
\[
F_N(r; k, \ell) = \frac{1}{\ell} + O_{k, \ell}(N^{-1/2}),
\]
\emph{unconditionally on H3$'$, requiring only H1$'$.}
\end{theorem}
\end{rigorousbox}

\begin{proof}
If $\ell \mid 2k$, then $2kq \equiv 0 \pmod \ell$ for all $q$, so
$p \equiv km(m+1) + \eps \pmod \ell$. We claim $km(m+1) \equiv 0 \pmod \ell$:
indeed, either $\ell \mid k$ (then trivially), or $\ell \mid 2$ giving $\ell = 2$,
in which case $m(m+1)$ is always even. In either case, $p \equiv \eps \neq 0 \pmod \ell$
since $\gcd(\eps, 2k) = 1$. The uniform distribution then follows from H1$'$ alone,
the $\ell - 1$ permitted residues becoming the $\ell$ residues, all available.
\end{proof}

\begin{remark}
For $k = 3$, Theorem~\ref{thm:degenerate} applies to $\ell \in \{2, 3\}$. For
$k = 21$, it applies to $\ell \in \{2, 3, 7\}$. For $k = 15$, to $\ell \in \{2, 3, 5\}$.
The empirical validation in Section~\ref{sec:universality} confirms the theorem
across all 10 cases tested.
\end{remark}

\begin{figure}[h!]
\centering
\includegraphics[width=0.92\textwidth]{fig8_two_regimes.png}
\caption{The two regimes of $q_{\min}$ modulo $\ell$ for $k = 3$. \textbf{(A)}
Degenerate case $\ell = 2 \mid 2k$: distribution exactly uniform
(Theorem~\ref{thm:degenerate}, rigorous, no hypothesis). \textbf{(B)} Non-degenerate
case $\ell = 5$: distribution non-uniform, predicted by
Theorem~\ref{thm:main} (heuristic, conditional on H1$'$+H3$'$).}
\label{fig:two-regimes}
\end{figure}

\subsection{Conjectural tighter bound}

\begin{openproblem}
\begin{conjecture}\label{conj:tight}
The effective constant in Theorem~\ref{thm:main} satisfies
\[
\limsup_{N \to \infty} \sqrt{N} \cdot \max_r \left|F_N(r) - \frac{\PiP(r)}{\ell - 1}\right|
\le C_{\mathrm{opt}}(k, \ell)
\]
with $C_{\mathrm{opt}}(k, \ell) \approx 0.5$, \emph{independent of $k$}, and roughly
$5\times$ smaller than $C(k, \ell)$ from Corollary~\ref{cor:constant}. The factor
$5$ originates from the constraint $\sum_{m_0, \eps} (\nu_N(m_0, \eps)
- 1/\ell|\Eclass_k|) = 0$, which the triangle inequality (Step~4 of the proof)
does not exploit.
\end{conjecture}
\textbf{Status.} Empirically validated at $N = 10^6$ for $k \in \{3, 7, 15, 21\}$
and $\ell \le 31$: average effective constants $\bar C_3 = 0.47$, $\bar C_7 = 0.51$,
$\bar C_{15} = 0.49$, $\bar C_{21} = 0.46$. Proof would require a Cauchy--Schwarz
refinement.
\end{openproblem}

\section{The Anti-Persistence Sign-Reversal Phenomenon}
\label{sec:anti-persistence}

This section reports a phenomenon distinct from the marginal $p$-adic structure
of Section~\ref{sec:framework}. While Theorems~\ref{thm:main} and~\ref{thm:degenerate}
concern the \emph{static} distribution of $\qmin \bmod \ell$, the present
phenomenon concerns the \emph{temporal} sequence $(\qmin)_n$ of the parametric
family ordered by increasing $p$.

\subsection{Empirical observation across 10 values of $k$}

We extended the validation to $k \in \{3, 5, 7, 9, 11, 13, 15, 17, 19, 21\}$
(all squarefree odd $k \le 21$ plus $k = 15$), each at $N = 10^6$.

\begin{discoverybox}
\begin{observation}[Linear law for $\mathrm{ACF}_1(q \bmod 2)$]\label{disc:acf}
At $N = 10^6$, the lag-1 autocorrelation of $(q_n \bmod 2)$ varies
\textbf{linearly with $k$}, not with $\varphi(2k)$:
\[
\mathrm{ACF}_1(q \bmod 2; k) \approx 0.0237\,k - 0.199 \qquad (R^2 = 0.97).
\]
The sign-reversal threshold is $k_0 \approx 8.4$: \emph{anti-persistence} for
$k < k_0$, \emph{positive persistence} for $k > k_0$. Empirically, $k=9$ is
nearly at the threshold ($\mathrm{ACF}_1 = -0.0075$).
\end{observation}
\end{discoverybox}

\begin{table}[h!]
\centering
\caption{Lag-1 autocorrelation of $(q \bmod 2)$ for $k \in \{3, 5, \ldots, 21\}$
at $N = 10^6$. Sign reversal occurs between $k = 9$ and $k = 11$.}
\begin{tabular}{c c c c c}
\toprule
$k$ & $\varphi(2k)$ & $\mathrm{ACF}_1$ & $\mathbb{P}(\text{same parity})$ & $z$-score (runs) \\
\midrule
3 & 2 & $-0.0719$ & $46.4\%$ & $+71.9\sigma$ \\
5 & 4 & $-0.0877$ & $45.6\%$ & $+87.7\sigma$ \\
7 & 6 & $-0.0572$ & $47.1\%$ & $+57.2\sigma$ \\
9 & 6 & $\mathbf{-0.0075}$ & $\mathbf{49.6\%}$ & $\mathbf{+7.5\sigma}$ \\
11 & 10 & $\mathbf{+0.0399}$ & $\mathbf{52.0\%}$ & $\mathbf{-39.9\sigma}$ \\
13 & 12 & $+0.0930$ & $54.7\%$ & $-93.0\sigma$ \\
15 & 8 & $+0.1574$ & $57.9\%$ & $-181\sigma$ \\
17 & 16 & $+0.2103$ & $60.5\%$ & $-216\sigma$ \\
19 & 18 & $+0.2643$ & $63.2\%$ & $-273\sigma$ \\
21 & 12 & $+0.3152$ & $65.8\%$ & $-323\sigma$ \\
\bottomrule
\end{tabular}
\end{table}

\begin{figure}[h!]
\centering
\includegraphics[width=0.97\textwidth]{fig9_acf_extended.png}
\caption{ACF$_1$ extended investigation (10 values of $k$, $N = 10^6$).
\textbf{(A)} Plotted vs.\ $k$, the data follow a remarkably clean linear law
$\mathrm{ACF}_1 \approx 0.024\,k - 0.20$ with $R^2 = 0.97$. The transition to
positive persistence occurs at $k \approx 8.4$. \textbf{(B)} Plotted vs.\
$\varphi(2k)$, the law is poor ($R^2 = 0.72$): the points $k = 13$ and $k = 21$
both have $\varphi(2k) = 12$ but $\mathrm{ACF}_1$ values differing by $0.22$.
This rules out $\varphi(2k)$ as the controlling variable.}
\label{fig:acf-extended}
\end{figure}

\subsection{The controlling variable is $k$, not $\varphi(2k)$}

A natural initial hypothesis was that $\mathrm{ACF}_1$ depends on
$|\Eclass_k| = \varphi(2k)$, the number of admissible $\eps$-classes. However,
the data refute this:
\begin{itemize}[leftmargin=*]
\item $k = 13$ and $k = 21$ both have $\varphi(2k) = 12$, yet
$\mathrm{ACF}_1(k=13) = +0.093$ vs.\ $\mathrm{ACF}_1(k=21) = +0.315$ --- a
difference of $0.22$.
\item $k = 7$ and $k = 9$ both have $\varphi(2k) = 6$, but
$\mathrm{ACF}_1(k=7) = -0.057$ vs.\ $\mathrm{ACF}_1(k=9) = -0.008$.
\item Linear regression on $\varphi(2k)$ alone gives $R^2 = 0.72$, while
linear regression on $k$ alone gives $R^2 = 0.97$.
\end{itemize}

The bivariate regression $\mathrm{ACF}_1 = a + b\,k + c\,\varphi(2k)$ confirms:
$b = +0.0249$ (significant), $c = -0.0015$ (negligible).

\subsection{Magnitude of effect}

For $k = 21$: over $10^6$ pairs $(q_n, q_{n+1})$, the observed number of
same-parity transitions is $657{,}600$, versus $500{,}000$ expected under
independence. This represents an excess of $157{,}600$, with Wald--Wolfowitz
$z$-score of $-323\sigma$. The phenomenon is not a small-$N$ artefact.

\subsection{Damped oscillation across lags: spectral signature}

Examining all autocorrelations from lag~1 to lag~10, a remarkable spectral
structure emerges. For each lag $\ell$, the linear regression $\mathrm{ACF}_\ell
\approx \alpha(\ell) \cdot k + \beta(\ell)$ gives slopes:

\begin{table}[h!]
\centering
\caption{Slopes $\alpha(\ell)$ of the linear regression of $\mathrm{ACF}_\ell$
on $k$, for $\ell = 1$ to $10$ (data from 10 values of $k$ at $N = 10^6$).}
\begin{tabular}{c c c c}
\toprule
lag $\ell$ & $\alpha(\ell)$ & $\beta(\ell)$ & $R^2$ \\
\midrule
1 & $\mathbf{+0.0237}$ & $-0.199$ & $\mathbf{0.971}$ \\
2 & $-0.0106$ & $-0.010$ & $0.717$ \\
3 & $\mathbf{-0.0179}$ & $+0.116$ & $\mathbf{0.839}$ \\
4 & $-0.0063$ & $+0.061$ & $0.396$ \\
5 & $+0.0041$ & $-0.023$ & $0.664$ \\
6 & $+0.0058$ & $-0.046$ & $0.661$ \\
7 & $+0.0025$ & $-0.026$ & $0.356$ \\
8 & $-0.0007$ & $-0.000$ & $0.225$ \\
9 & $-0.0013$ & $+0.006$ & $0.380$ \\
10 & $-0.0004$ & $-0.000$ & $0.075$ \\
\bottomrule
\end{tabular}
\end{table}

\begin{figure}[h!]
\centering
\includegraphics[width=0.97\textwidth]{fig11_acf_spectrum.png}
\caption{Spectral structure of $(q_n \bmod 2)$ autocorrelation. \textbf{(A)}
The slopes $\alpha(\ell) = \partial \mathrm{ACF}_\ell / \partial k$ exhibit
the pattern $(+, -, -, -, +, +, +, 0, 0, 0)$, characteristic of a damped
oscillation. Stars indicate significance: $***$ for $R^2 > 0.8$,
$**$ for $R^2 > 0.5$, $*$ for $R^2 > 0.3$. \textbf{(B)} The full ACF spectrum
for $k \in \{15, 17, 21\}$ shows the same shape but with amplitude scaling
linearly in $k$.}
\label{fig:acf-spectrum}
\end{figure}

\begin{discoverybox}
\begin{observation}[Damped cosine signature]\label{disc:cosine}
The slopes $\alpha(\ell)$ are remarkably well fitted by a damped cosine:
\[
\alpha(\ell) \approx 0.039 \cdot e^{-0.39(\ell-1)} \cdot \cos\!\bigl(1.07(\ell-1) + 0.92\bigr),
\qquad R^2 = 0.9988.
\]
The fit reveals a \textbf{period} $T = 2\pi/\omega \approx 5.89$ and a
\textbf{damping rate} $\gamma \approx 0.39$ per lag. This is a remarkably
clean spectral signature, suggesting an underlying analytic mechanism.
\end{observation}
\end{discoverybox}

\begin{figure}[h!]
\centering
\includegraphics[width=0.85\textwidth]{fig13_damped_cosine.png}
\caption{The slopes $\alpha(\ell)$ are fitted to within $R^2 = 0.9988$ by
a damped cosine of period $T \approx 5.89$ and damping rate $\gamma \approx 0.39$.
Each empirical point sits on the fitted curve. This level of precision over
10 data points is highly non-trivial.}
\label{fig:damped-cosine}
\end{figure}

\subsection{Amplification with modulus $\ell$}

Repeating the analysis modulo $\ell \in \{3, 5\}$ reveals that the lag-1
autocorrelation \emph{amplifies} dramatically with $\ell$:

\begin{table}[h!]
\centering
\caption{Lag-1 autocorrelation of $(q \bmod \ell)$ for $\ell \in \{2, 3, 5\}$
across selected $k$. Cells marked ``--'' correspond to degenerate cases
($\ell \mid 2k$).}
\begin{tabular}{c c c c}
\toprule
$k$ & $\mathrm{ACF}_1(\bmod 2)$ & $\mathrm{ACF}_1(\bmod 3)$ & $\mathrm{ACF}_1(\bmod 5)$ \\
\midrule
3 & $-0.072$ & -- & $-0.062$ \\
7 & $-0.057$ & $+0.007$ & $+0.165$ \\
11 & $+0.040$ & $+0.158$ & $+0.376$ \\
13 & $+0.093$ & $+0.239$ & $+0.448$ \\
17 & $+0.210$ & $+0.362$ & $+0.559$ \\
19 & $+0.264$ & $+0.422$ & $+0.600$ \\
21 & $\mathbf{+0.315}$ & -- & $\mathbf{+0.643}$ \\
\bottomrule
\end{tabular}
\end{table}

\begin{figure}[h!]
\centering
\includegraphics[width=0.78\textwidth]{fig12_acf_modulus.png}
\caption{Amplification of $\mathrm{ACF}_1$ with the modulus $\ell$. The three
curves (mod 2, 3, 5) follow the same general linear-in-$k$ shape, but with
amplitude growing in $\ell$. For $k = 21$ at $\ell = 5$, $\mathrm{ACF}_1
\approx +0.64$, twice the mod-2 value.}
\label{fig:acf-modulus}
\end{figure}

For $k = 21$, mod 5: $\mathrm{ACF}_1 = +0.64$, meaning that
$\mathbb{P}(q_n \equiv q_{n+1} \pmod 5)$ exceeds $0.32$, double the
independence value $0.20$. This is one of the strongest deterministic effects
observed in the parametric prime family at this scale.

\subsection{Refined conjecture (extended)}

\begin{openproblem}
\begin{conjecture}[Spectral structure of $(q_n \bmod \ell)$]\label{conj:acf}
For the parametric prime family with $\ell \in \{2, 3, 5, \ldots\}$ and
$\gcd(\ell, 2k) = 1$, as $N \to \infty$:
\begin{enumerate}[leftmargin=*]
\item \textbf{Linear law per lag.} The autocorrelation at lag $j$ satisfies
$\mathrm{ACF}_j(q \bmod 2; k) \approx \alpha_2(j)\,k + \beta_2(j)$, with
$\alpha_2(1) = +0.024$, $\beta_2(1) = -0.20$.
\item \textbf{Damped cosine spectrum.} The slope sequence $\{\alpha_2(j)\}_{j=1}^{10}$
is fitted by
\[
\alpha_2(j) \approx A \cdot e^{-\gamma(j-1)} \cdot \cos\bigl(\omega(j-1) + \varphi\bigr),
\quad A \approx 0.039,\ \gamma \approx 0.39,\ \omega \approx 1.07,\ \varphi \approx 0.92,
\]
to within $R^2 = 0.9988$. The period $T = 2\pi/\omega \approx 5.89$.
\item \textbf{Modulus amplification.} For each $k$ admissible,
$\mathrm{ACF}_1(q \bmod 3; k) > \mathrm{ACF}_1(q \bmod 2; k)$ when $\gcd(3, 2k) = 1$;
similarly for mod 5. The amplification factor is approximately $2\times$
between mod 2 and mod 5 at large $k$.
\end{enumerate}
\end{conjecture}
\end{openproblem}

\textbf{Status.} Established empirically on 10 values of $k \in [3, 21]$.
The damped cosine fit is suggestive of an analytic origin (\textit{e.g.}, a
specific spectral decomposition of the joint distribution of consecutive
$\qmin$ values), and is the most striking quantitative finding of this section.

\section{Empirical Validation: $k = 3$ at $N = 10^6$}
\label{sec:validation-k3}

\subsection{Database}

A segmented Eratosthenes sieve generated all primes $p \le 1.55 \times 10^7$,
yielding $10^6$ triples $(m, \eps, q)$ for $k = 3$ in 1.8 s. Statistics:
$m_{\max} = 2{,}271$, $|q|_{\max} = 1{,}135$, $|q|_{\mathrm{mean}} = 369.7$.
This range covers values of $|q|$ much larger than any tested modulus $\ell$,
exiting the window-truncation regime.

\subsection{Validation of Theorem~\ref{thm:main} (non-degenerate case)}

\begin{table}[h!]
\centering
\caption{Validation at $N = 10^6$ for $k = 3$. Pearson correlation $r$ between
predicted $\PiP(r)/(\ell-1)$ and observed $F_N(r)$, and effective constant
$C_{\mathrm{eff}} = \mathrm{err}_{\max} \cdot \sqrt{N}$.}
\begin{tabular}{c c c c c}
\toprule
$\ell$ & Pearson $r$ & err$_{\max}$ & $C_{\mathrm{eff}}$ & Bound $C(3, \ell)$ \\
\midrule
5 & 0.99996 & $2.96 \times 10^{-4}$ & 0.30 & 2.37 \\
7 & 0.99992 & $6.78 \times 10^{-4}$ & 0.68 & 2.25 \\
11 & 0.99917 & $5.16 \times 10^{-4}$ & 0.52 & 2.15 \\
13 & 0.99921 & $4.40 \times 10^{-4}$ & 0.44 & 2.12 \\
17 & 0.99572 & $5.49 \times 10^{-4}$ & 0.55 & 2.10 \\
19 & 0.99541 & $3.79 \times 10^{-4}$ & 0.38 & 2.09 \\
23 & 0.99182 & $3.92 \times 10^{-4}$ & 0.39 & 2.08 \\
\bottomrule
\end{tabular}
\end{table}

The Pearson correlation exceeds $0.991$ for all $\ell \le 23$, and the effective
constant remains stable around $\bar C \approx 0.5$, confirming
Conjecture~\ref{conj:tight}.

\begin{figure}[h!]
\centering
\includegraphics[width=0.98\textwidth]{fig1_distribution_k3.png}
\caption{Predicted vs.\ observed distribution of $q_{\min} \bmod \ell$ for
$k = 3$, $N = 10^6$, and $\ell \in \{5, 7, 11, 13\}$. The Pearson correlation
exceeds $0.999$ in all cases. The dashed line marks the naive uniform value
$1/\ell$.}
\label{fig:dist-k3}
\end{figure}

\subsection{Detail at $k = 3$, $\ell = 5$}

\begin{table}[h!]
\centering
\caption{Predicted vs.\ observed at $k = 3$, $\ell = 5$, $N = 10^6$.}
\begin{tabular}{c c c c}
\toprule
$r$ & $\PiP(r)/(\ell-1)$ predicted & $F_N(r)$ observed & Error \\
\midrule
0 & 0.200000 & 0.199761 & $2.4 \times 10^{-4}$ \\
1 & 0.175000 & 0.175296 & $3.0 \times 10^{-4}$ \\
2 & 0.250000 & 0.250289 & $2.9 \times 10^{-4}$ \\
3 & 0.175000 & 0.174906 & $9.4 \times 10^{-5}$ \\
4 & 0.200000 & 0.199748 & $2.5 \times 10^{-4}$ \\
\bottomrule
\end{tabular}
\end{table}

The agreement is to $4$ decimal places, with maximal error
$2.96 \times 10^{-4} \approx 0.3 / \sqrt{N}$.

\subsection{Scaling test in $N$}

\begin{table}[h!]
\centering
\caption{Scaling of $\mathrm{err}_{\max}$ vs.\ $N$ for $k=3$, $\ell=5$. The product
$\mathrm{err}_{\max} \cdot \sqrt{N}$ should remain bounded.}
\begin{tabular}{c c c}
\toprule
$N$ & err$_{\max}$ & err$_{\max} \cdot \sqrt{N}$ \\
\midrule
$10^2$ & $4.50 \times 10^{-2}$ & 0.450 \\
$10^3$ & $8.00 \times 10^{-3}$ & 0.253 \\
$10^4$ & $3.00 \times 10^{-3}$ & 0.300 \\
$10^5$ & $1.06 \times 10^{-3}$ & 0.335 \\
$10^6$ & $2.96 \times 10^{-4}$ & 0.296 \\
\bottomrule
\end{tabular}
\end{table}

The product remains in the band $[0.25, 0.45]$ across four decades of $N$,
robustly confirming the $N^{-1/2}$ rate.

\begin{figure}[h!]
\centering
\includegraphics[width=0.95\textwidth]{fig2_scaling.png}
\caption{Scaling of err$_{\max}$ with $N$ for $k = 3$, $\ell = 5$.
\textbf{(A)} Log-log plot: err$_{\max}$ closely follows the $0.3 / \sqrt{N}$
heuristic; the Cor.~\ref{cor:constant} bound $C(3, 5)/\sqrt{N} = 2.37/\sqrt{N}$
is much larger. \textbf{(B)} The product $\mathrm{err}_{\max} \cdot \sqrt{N}$
remains stable in the band $[0.25, 0.45]$ across four decades of $N$, confirming
the $N^{-1/2}$ rate of Theorem~\ref{thm:main}.}
\label{fig:scaling}
\end{figure}

\subsection{Validation of H3$'$ at the stratum level}

For the 10 strata $(m_0, \eps) \in \Z/5\Z \times \{\pm 1\}$ at $N = 10^6$: total
violations $= 1$ (one $\qmin$ landing on the forbidden residue out of $10^6$
cases); all 10 strata pass the $\chi^2$ uniformity test on the four permitted
residues with $p \ge 0.61$. Hypothesis~\ref{hyp:H3prime} thus holds to $99.9999\%$
empirically.

\begin{figure}[h!]
\centering
\includegraphics[width=0.98\textwidth]{fig6_h3prime.png}
\caption{Empirical validation of Hypothesis~H3$'$ for $k = 3$, $\ell = 5$, $N = 10^6$.
For each of the 10 strata $(m_0, \eps) \in \Z/5\Z \times \{\pm 1\}$, the
forbidden residue $q^*$ (red) is empty (1 violation total out of $10^6$ cases),
and the four permitted residues (blue) are uniformly distributed (dashed line),
with $\chi^2$ $p$-values $\ge 0.61$ in every stratum.}
\label{fig:h3prime}
\end{figure}

\section{Universality across $k \in \{3, 7, 15, 21\}$}
\label{sec:universality}

\subsection{Database statistics}

\begin{table}[h!]
\centering
\caption{Database statistics at $N = 10^6$ for each $k$.}
\begin{tabular}{c c c c c c}
\toprule
$k$ & $|\Eclass_k| = \varphi(2k)$ & $m_{\max}$ & $|q|_{\max}$ & $|q|_{\mathrm{mean}}$ & Generation time \\
\midrule
3 & 2 & 2{,}271 & 1{,}135 & 369.7 & 1.8 s \\
7 & 6 & 1{,}487 & 743 & 242.1 & 2.8 s \\
15 & 8 & 1{,}016 & 508 & 165.3 & 2.4 s \\
21 & 12 & 858 & 429 & 139.7 & 2.4 s \\
\bottomrule
\end{tabular}
\end{table}

\subsection{Validation of Theorem~\ref{thm:degenerate} (degenerate case)}

\begin{table}[h!]
\centering
\caption{Theorem~\ref{thm:degenerate} validation: 10 cases tested, all pass.}
\begin{tabular}{c c c c c c}
\toprule
$k$ & $\ell \mid 2k$ & $\chi^2$ & $p$-value & Max bias (\%) & Verdict \\
\midrule
3 & 2 & 0.34 & 0.562 & $+0.058$ & uniform \checkmark \\
3 & 3 & 0.76 & 0.685 & $+0.115$ & uniform \checkmark \\
7 & 2 & 3.07 & 0.080 & $+0.175$ & uniform \checkmark \\
7 & 7 & 1.47 & 0.962 & $+0.157$ & uniform \checkmark \\
15 & 2 & 2.46 & 0.116 & $+0.157$ & uniform \checkmark \\
15 & 3 & 0.21 & 0.899 & $+0.036$ & uniform \checkmark \\
15 & 5 & 0.38 & 0.984 & $+0.097$ & uniform \checkmark \\
21 & 2 & 0.02 & 0.897 & $+0.013$ & uniform \checkmark \\
21 & 3 & 0.19 & 0.911 & $+0.053$ & uniform \checkmark \\
21 & 7 & 2.78 & 0.836 & $+0.264$ & uniform \checkmark \\
\bottomrule
\end{tabular}
\end{table}

\begin{figure}[h!]
\centering
\includegraphics[width=0.98\textwidth]{fig5_degenerate.png}
\caption{Validation of Theorem~\ref{thm:degenerate} (degenerate case $\ell \mid 2k$)
across all 10 cases tested. In each panel, the predicted uniform distribution
(blue, $1/\ell$) and the observed distribution (orange) are indistinguishable;
maximum bias $\le 0.27\%$.}
\label{fig:degenerate}
\end{figure}

\subsection{Validation of Theorem~\ref{thm:main} (universality)}

\begin{table}[h!]
\centering
\caption{Pearson correlation $r$(predicted, observed) for $k \in \{3, 7, 15, 21\}$
at $N = 10^6$. \textbf{*} = degenerate case (Theorem~\ref{thm:degenerate} applies).}
\begin{tabular}{c c c c c}
\toprule
$\ell$ & $k=3$ & $k=7$ & $k=15$ & $k=21$ \\
\midrule
5 & 0.99996 & 0.99968 & * & 0.99844 \\
7 & 0.99992 & * & 0.99933 & * \\
11 & 0.99917 & 0.99609 & 0.99628 & 0.98470 \\
13 & 0.99921 & 0.98866 & 0.99450 & 0.98179 \\
17 & 0.99572 & 0.98282 & 0.95749 & 0.95810 \\
19 & 0.99541 & 0.96411 & 0.96382 & 0.94207 \\
23 & 0.99182 & 0.95754 & 0.92746 & 0.92597 \\
29 & 0.96673 & 0.92006 & 0.89323 & 0.77188 \\
31 & 0.96103 & 0.83416 & 0.90752 & 0.70680 \\
\midrule
\textbf{Mean} & \textbf{0.990} & \textbf{0.955} & \textbf{0.955} & \textbf{0.909} \\
\bottomrule
\end{tabular}
\end{table}

\paragraph{Observation.} The correlation is excellent ($r > 0.99$) for all
$\ell \le 13$ and all $k$. Degradation appears for large $(\ell, k)$, consistent
with the theoretical constant $C(k, \ell) = \ell |\Eclass_k|/(\ell-1) \cdot \sigma$
which grows with $k$ and $\ell$. To validate Theorem~\ref{thm:main} for $k=21$,
$\ell \ge 29$, larger $N$ ($10^7$ or more) would be required (Open Problem~C).

\begin{figure}[h!]
\centering
\includegraphics[width=0.98\textwidth]{fig3_universalite.png}
\caption{Universality test of Theorem~\ref{thm:main} across $k \in \{3, 7, 15, 21\}$.
\textbf{(A)} Pearson correlation $r$(predicted, observed) vs.\ $\ell$. For $\ell \le 13$,
$r > 0.99$ for all $k$. Degradation appears for $\ell \ge 17$ and large $k$, consistent
with the growth of $C(k, \ell)$. \textbf{(B)} Effective constant
$C_{\mathrm{eff}} = \mathrm{err}_{\max} \cdot \sqrt{N}$ vs.\ $\ell$. Despite fluctuations,
$C_{\mathrm{eff}}$ stays close to the conjectured optimal $\bar C \approx 0.5$
(Conjecture~\ref{conj:tight}), independent of $k$.}
\label{fig:universalite}
\end{figure}

\begin{figure}[h!]
\centering
\includegraphics[width=0.98\textwidth]{fig7_predicted_observed.png}
\caption{Predicted (blue) vs.\ observed (orange) distribution of $q_{\min} \bmod \ell$
for various $(k, \ell)$ at $N = 10^6$. Top row: $k = 3$, $\ell \in \{5, 7, 11, 13\}$.
Bottom row: $k = 21$, $\ell \in \{5, 11, 17, 23\}$. The dotted line marks the naive
uniform value $1/\ell$. The non-trivial structure of $\pi_{\mathrm{perm}}$ is
visible in every panel; the agreement between prediction and observation is excellent
for $\ell \le 13$ and remains good for larger $\ell$.}
\label{fig:predicted-observed}
\end{figure}

\subsection{Detail at $k = 21$, $\ell = 5$}

\begin{table}[h!]
\centering
\caption{Predicted vs.\ observed at $k = 21$, $\ell = 5$, $N = 10^6$.}
\begin{tabular}{c c c c}
\toprule
$r$ & $\PiP(r)/(\ell-1)$ & $F_N(r)$ & Error \\
\midrule
0 & 0.204167 & 0.203777 & $3.9 \times 10^{-4}$ \\
1 & 0.200000 & 0.199898 & $1.0 \times 10^{-4}$ \\
2 & 0.191667 & 0.191745 & $7.8 \times 10^{-5}$ \\
3 & 0.200000 & 0.200018 & $1.8 \times 10^{-5}$ \\
4 & 0.204167 & 0.204562 & $4.0 \times 10^{-4}$ \\
\bottomrule
\end{tabular}
\end{table}

The non-uniform predicted distribution (residue~2 least likely, residues~0
and~4 most likely) is matched to 4 decimal places, with maximal error
$3.95 \times 10^{-4}$.

\section{Connection to Paper~II's Hypothesis H3}
\label{sec:connection}

The aim of this section is to establish that Hypothesis~\ref{hyp:H3prime} (H3$'$)
is a \emph{consequence} of Hypothesis~H3 of Paper~II, modulo the standard
Bateman--Horn machinery. The argument is heuristic but quantitative: under
H3 plus Cramér--Granville-type independence (H2), H3$'$ holds with explicit
error bounds. The key consequence is that Theorem~\ref{thm:main} provides
an \emph{indirect quantitative test} of H3.

\subsection{Statement of the implication}

\begin{heuristicbox}
\begin{proposition}[H3 implies H3$'$ at leading order]\label{prop:h3-implies-h3prime}
Let $k \ge 1$ and $\ell$ a prime with $\gcd(2k, \ell) = 1$. Assume H1 and H2
(Paper~I) and H3 (Paper~II) hold. Then Hypothesis~H3$'$
(Hypothesis~\ref{hyp:H3prime}) holds, with error term:
\[
\left| \Prb\!\left(\qmin \equiv r \pmod \ell \mid m \equiv m_0, \eps\right)
- \frac{\mathbb{1}[r \neq q^*(k, m_0, \eps; \ell)]}{\ell - 1} \right|
= O_{k, \ell}\!\left(\frac{1}{m \log m}\right).
\]
Aggregated over $N$ primes of the family, this contributes an error of order
$O(N^{-1/2} / \log N)$ to the bound of Theorem~\ref{thm:main}, strictly
sub-dominant relative to the H1$'$ statistical error of $O(N^{-1/2})$.
\end{proposition}
\end{heuristicbox}

\subsection{Proof outline (5 steps)}

The proof proceeds in five steps, paralleling the structure of Step~4 in the
proof of Theorem~\ref{thm:main}.

\paragraph{Step 1 (Reduction to densities in arithmetic progressions).}
For fixed $(m, \eps)$ and a permitted residue $r \neq q^*(k, m_0, \eps; \ell)$,
the candidates $q$ satisfying $q \equiv r \pmod \ell$ form an arithmetic
progression $r + \ell \Z$. The corresponding values of $p$ form the polynomial
sequence in $j$:
\[
g_{r, \eps}(j) = H_m + \eps + 2k(r + \ell j), \qquad j \in \Z,
\]
where $H_m = km(m+1)$. The probability that $\qmin \equiv r$ equals the
probability that the smallest $|q|$ giving a prime $p_{k,m,\eps,q}$ has
$q \equiv r \pmod \ell$. Under independence of primality events (H2), this
reduces, by the standard geometric distribution argument, to the relative
densities of primes in the $\ell-1$ permitted classes.

\paragraph{Step 2 (Bateman--Horn for the shifted polynomial).}
Apply Bateman--Horn (H1) to $g_{r, \eps}(j) = (2k\ell)\,j + (H_m + \eps + 2kr)$.
This is a degree-1 polynomial in $j$ with linear coefficient $2k\ell$ and
constant term $A_{r, \eps} := H_m + \eps + 2kr$. The Bateman--Horn singular
series is
\[
C(g_{r, \eps}) = \prod_{p \text{ prime}}
\frac{p - \omega_{g_{r,\eps}}(p)}{p - 1},
\]
where $\omega_{g_{r, \eps}}(p) = \#\{j \in \Z/p\Z : g_{r, \eps}(j) \equiv 0 \pmod p\}$.

\paragraph{Step 3 (Computation of local factors).}
We compute $\omega_{g_{r, \eps}}(p)$ explicitly:
\begin{itemize}[leftmargin=*]
\item For $p \mid 2k\ell$: case-by-case analysis using $\gcd(\eps, 2k) = 1$.
The contributions are independent of $r$ since they depend only on
$A_{r, \eps} \bmod p$ for $p \mid 2k\ell$.
\item For $p \nmid 2k\ell$: the equation $(2k\ell) j \equiv -A_{r, \eps} \pmod p$
has a \emph{unique} solution since $\gcd(2k\ell, p) = 1$, hence
$\omega_{g_{r, \eps}}(p) = 1$. The local factor is $(p-1)/(p-1) = 1$,
\emph{independent of $r$}.
\item For $p = \ell$: $g_{r, \eps}(j) \equiv H_m + \eps + 2kr \pmod \ell$,
\emph{constant in $j$}. By definition of $q^*$, this constant is non-zero
$\bmod \ell$ iff $r \neq q^*$. For permitted $r$, $\omega_{g_{r, \eps}}(\ell) = 0$
and the local factor is $\ell/(\ell-1)$, \emph{equal across all permitted $r$}.
\end{itemize}

\paragraph{Step 4 (Application of H3).}
The local factor at $\ell$ is the same for all permitted $r$ (Step~3). Therefore,
\emph{the only source of $r$-dependence in $C(g_{r, \eps})$ comes from primes
$p \nmid 2k\ell$ that contribute via the constant term $A_{r, \eps}$}.
This is precisely the situation governed by Hypothesis~H3 of Paper~II, applied
to the shifted polynomials $g_{r, \eps}$. H3 asserts that:
\[
\left| C(g_{r, \eps}) - C(g_{r', \eps}) \right|
= O\!\left(\frac{1}{m \log m}\right) \cdot C(g_{r, \eps})
\]
for any pair of permitted residues $r, r'$.

\paragraph{Step 5 (Conditional probability).}
Combining Steps 1--4: the densities of primes in the $(\ell-1)$ permitted
classes are equal up to multiplicative error $1 + O(1/(m \log m))$. By the
geometric argument of Step~1, this propagates to:
\[
\Prb\!\left(\qmin \equiv r \pmod \ell \mid m \equiv m_0, \eps\right)
= \frac{1}{\ell - 1} + O_{k, \ell}\!\left(\frac{1}{m \log m}\right),
\]
which is precisely the content of H3$'$. \qed

\subsection{Aggregation and propagation of error}

When $F_N(r; k, \ell)$ is computed via Theorem~\ref{thm:main}, the error from
H3 propagates as follows. The maximal layer index for the first $N$ primes
satisfies $M(N) \sim \sqrt{N \log N / k}$ (Conjecture~C3 of Paper~I). The H3
error per layer is $O(1/(M \log M))$, and the average over layers contributes
$O(1/(M \log M)) = O(1/(\sqrt{N \log N} \cdot \log\sqrt{N})) = O(1/(\sqrt{N}
(\log N)^{3/2}))$, which is strictly sub-dominant to the H1$'$ statistical
error $O(N^{-1/2})$.

\subsection{Empirical confirmation of the reduction}

The proof relies on the equality (modulo H3) of prime densities in the $\ell-1$
permitted residue classes. We verified this directly for $k = 3$, $\ell = 5$:

\begin{table}[h!]
\centering
\caption{Density of primes per permitted residue class, fixed $(m, \eps)$,
sampled at 100 values of $m \in [50, 5050]$.}
\begin{tabular}{l c c}
\toprule
Quantity & Value & Statistical prediction \\
\midrule
Mean ratio $n_{r=0} / n_{r=1}$ over 100 m & $1.0064$ & $1.000$ \\
Std.\ dev.\ of ratio & $0.0647$ & $1/\sqrt{200} = 0.0707$ \\
Bias from unity & $0.64\%$ & $0\%$ \\
\bottomrule
\end{tabular}
\end{table}

The mean ratio is unity to $0.64\%$, and the standard deviation matches the
purely statistical prediction $1/\sqrt{n_{\text{per class}}}$ to within $10\%$.
The H3 hypothesis holds at this level of precision empirically.

\subsection{Indirect quantitative test of H3}

The empirical verification of Theorem~\ref{thm:main} at $N = 10^6$
(Section~\ref{sec:validation-k3}: Pearson $r > 0.999$ for $\ell \le 13$,
$k = 3$) is consistent with H3$'$ holding in the regime tested, and (by
Proposition~\ref{prop:h3-implies-h3prime}) with H3 holding. The contrapositive
is the operational statement: a failure of Theorem~\ref{thm:main} at large
$N$ would imply a failure of H3$'$, which would in turn force a revision of
H3.

We deliberately avoid stronger language. As discussed in
Section~\ref{sec:second-order-warning}, no finite-$N$ numerical experiment
can settle a Bateman--Horn-type conjecture; the historical record contains
multiple cases where empirical patterns valid up to enormous $N$ failed
asymptotically (Skewes 1933 on $\pi(x) - \mathrm{Li}(x)$; the original
Mertens conjecture; certain Hardy--Littlewood prime-pair conjectures
shown to be mutually incompatible). Our results add empirical weight to
H3 in a specific, quantifiable manner; they neither prove H3 nor preclude
its eventual unconditional proof or disproof.

\subsection{The unconditional barrier}

The proof of Proposition~\ref{prop:h3-implies-h3prime} is heuristic at three
points: (i) it uses Bateman--Horn (H1), (ii) it uses Cramér--Granville
independence (H2), (iii) it uses H3. None of these is unconditionally proved,
and we wish to be precise about \emph{why} an unconditional version is currently
out of reach. Three obstacles compound:

\paragraph{(a) Bateman--Horn itself is open.} The only fully proved case of the
Bateman--Horn conjecture is the linear case $f(n) = an + b$
(Dirichlet~\cite{dirichlet} and the prime number theorem in arithmetic
progressions). Already $f(n) = n^2 + 1$ is open, despite over a century of
attention. Theorem~\ref{thm:degenerate} (the degenerate case of this paper) is
rigorous \emph{precisely because} it reduces to the linear case (Step~3 of the
proof of Proposition~\ref{prop:h3-implies-h3prime} shows that $g_{r,\eps}$ has
degree~1 in $j$); Theorem~\ref{thm:main} is heuristic precisely because it
needs Bateman--Horn for non-degenerate cases.

\paragraph{(b) The parity problem.} A more fundamental obstacle, identified by
Selberg~\cite{selberg} and made explicit in modern surveys
(\textit{e.g.}~\cite{friedlander, iwaniec}), is that sieve methods cannot
distinguish numbers with an odd number of prime factors from those with an
even number. This blocks any sieve-based proof of asymptotics for primes
represented by polynomials of degree $\ge 2$. The parity barrier is the
principal reason why Bateman--Horn---and a fortiori H3 of Paper~II and H3$'$
of the present paper---resist proof. Recent work (Maynard~\cite{maynard},
Friedlander--Iwaniec~\cite{friedlander}) has produced striking results
\emph{around} the barrier (small gaps, primes with given combinatorial
constraints), but has not broken it.

\paragraph{(c) Distribution in arithmetic progressions of large moduli.}
Bombieri--Vinogradov~\cite{bombieri, vinogradov-bv} provides PNT in arithmetic
progressions of moduli up to $N^{1/2-\epsilon}$ \emph{on average}. To prove H3
or H3$'$ unconditionally, one would need control \emph{uniformly} (not just
on average) over moduli growing with the layer index $m$, and---worse---in a
manner that interacts with the polynomial structure. This sits beyond the
reach of all currently known methods, including the Elliott--Halberstam
conjecture as a working hypothesis.

\subsection{The non-linearity of $\qmin$}

A subtler issue, specific to our setting, deserves emphasis. The quantity
$\qmin(k, m, \eps)$ is the \emph{minimum} of $|q|$ over all $q$ giving a prime;
it is therefore a non-linear functional of the primality indicator sequence
$(\mathbb{1}[p_{k,m,\eps,q} \in \Prim])_q$. Step~5 of the proof of
Proposition~\ref{prop:h3-implies-h3prime} appeals informally to the
``geometric distribution argument'': if the densities of primes in the
$\ell-1$ permitted classes are equal, then by symmetry the minimum-$|q|$ prime
is uniformly distributed over the permitted residues.

This argument is correct under \emph{exact independence} of primality events
across different $q$. Under realistic, finite-size dependence, it is at most
a leading-order approximation. The precise correction term is what
H3$'$ asserts to be sub-leading; but our argument does not produce a
quantitative bound on the dependence-induced second-order term. Such a bound
would require analytic control over correlations of the primality function on
short, polynomially-spaced sequences, which is itself an open problem.

\subsection{A second-order warning}
\label{sec:second-order-warning}

The most subtle danger is not that H3$'$ might be qualitatively wrong, but
that there might be a \emph{second-order} dependence not captured by H1+H2+H3.
This is a recurring failure mode in heuristic number theory: cf.\ the
Hardy--Littlewood conjectures, some pairs of which have been shown to be
mutually incompatible despite the considerable empirical support each
enjoyed in isolation. We point out that \emph{the empirically observed
spectral signature of Section~\ref{sec:anti-persistence}---the damped cosine
fit with $R^2 = 0.999$---may itself be the trace of a non-trivial second-order
correlation that H2 ignores}. Whether this reflects a genuine analytic
mechanism or merely an artefact of the averaging procedure is precisely the
content of Open Problem~D.

\subsection{Status: what numerical evidence does and does not establish}

The empirical correlation $r = 0.99996$ (Section~\ref{sec:validation-k3}) and
the $1/\sqrt{N}$ scaling over four decades constitute \textbf{strong but not
conclusive evidence} for Theorem~\ref{thm:main} and Hypothesis~H3$'$.
Numerical evidence cannot, by itself, settle a Bateman--Horn-type conjecture:
the historical record (Hardy--Littlewood incompatibilities, Skewes-type
sign-changes, Mertens conjecture) shows that asymptotic behaviour can deviate
from finite-$N$ patterns in ways invisible up to extremely large $N$. Our
results should therefore be read as: \emph{a precise quantitative version of
a Bateman--Horn-style heuristic, validated to high precision in the regime
$N \le 10^6$, with explicit error rates compatible with the natural conjectural
rates}. Their independent value is to provide testable consequences: any
future analytic progress on H3 (or its analogues) should reproduce the
predictions of Theorem~\ref{thm:main} as a special case.

\paragraph{What is, then, rigorously settled in this paper?}
The genuinely unconditional content is Theorem~\ref{thm:degenerate} (the
degenerate case $\ell \mid 2k$), which reduces to BV at fixed modulus and is
rigorous in the strict sense. Everything else---Theorem~\ref{thm:main},
Proposition~\ref{prop:h3-implies-h3prime}, Conjectures~\ref{conj:tight}
and~\ref{conj:acf}---is conditional or heuristic, with carefully tracked
hypotheses and explicit empirical validation.

\section{Negative Results}
\label{sec:negative}

\subsection{Window-truncation artefact for small $|q|$}

\begin{negativebox}
\begin{observation}[Window-truncation artefact, retracted]
At $N = 10^4$ for $k \ge 7$, the apparent ``U-shaped bias'' in the histogram
of $q \bmod 11$ is not a genuine $p$-adic bias but a geometric artefact: the
maximum $|q|$ at this scale is small (e.g., $|q|_{\max} = 13$ for $k = 21$,
$N = 10^4$), so the residue $q \bmod 11$ is essentially $q$ itself, distributed
in a bell shape around $0$ rather than uniformly across the residues.
\end{observation}
\end{negativebox}

\textbf{Resolution at $N = 10^6$:} the average $|q|_{\mathrm{mean}}$ rises to
$\approx 140$--$370$ depending on $k$, exiting the truncation regime. The
$p$-adic bias revealed in this regime matches Theorem~\ref{thm:main}.

\subsection{Spurious correlations from autocorrelated cumulative sequences}

Following the methodology of Paper~I (Section~7), we verified that the lag-1
ACF of $(q \bmod 2)$ is not a spurious effect of cumulative autocorrelation.
The Wald--Wolfowitz runs test, which counts transitions rather than cumulative
sums, gives $z > 11{,}000\sigma$ at $N = 10^6$ for $k = 21$, well beyond any
spurious-regression threshold.

\section{Open Problems}
\label{sec:openpb}

\begin{openproblem}
\textbf{Open Problem A.} Prove H3$'$ unconditionally for $k = 3$ and small primes
$\ell$. Equivalent in difficulty to Open Problem~12.2 of Paper~II
(Selberg sieve uniform in $(m_0, \eps)$, $N^{1/2}$ barrier).
\end{openproblem}

\begin{openproblem}
\textbf{Open Problem B.} Prove Conjecture~\ref{conj:tight} on the optimal
constant: $C_{\mathrm{opt}}(k, \ell) \approx 0.5$ independent of $k$. Equivalent
to controlling the $L^2$-norm of the deviation vector
$(\nu_N(m_0, \eps) - 1/(\ell |\Eclass_k|))_{(m_0, \eps)}$ under the linear
constraint $\sum = 1$, via Cauchy--Schwarz instead of triangle inequality.
\end{openproblem}

\begin{openproblem}
\textbf{Open Problem C.} Validate Theorem~\ref{thm:main} for $k = 21$ and
$\ell \ge 29$ at $N \ge 10^7$. The current degradation of $r$ to $0.71$ at
$\ell = 31$ is consistent with $1/\sqrt{N}$ scaling but requires direct
verification at larger scale.
\end{openproblem}

\begin{openproblem}
\textbf{Open Problem D (new, refined).} Derive Conjecture~\ref{conj:acf}
theoretically. Specifically: (i) prove the linear-in-$k$ law for
$\mathrm{ACF}_1(q \bmod 2)$; (ii) explain the damped cosine spectrum
$\alpha_2(\ell) \approx 0.039\,e^{-0.39(\ell-1)}\cos(1.07(\ell-1) + 0.92)$
with $R^2 = 0.999$, identifying the period $T \approx 5.89$ and damping rate
$\gamma \approx 0.39$ analytically; (iii) explain the modulus amplification
($\sim 2\times$ between mod 2 and mod 5). A natural starting point is the joint
distribution of $(\qmin)_n$ and $(\qmin)_{n+1}$ under Bateman--Horn for two
consecutive primes of the family.
\end{openproblem}

\begin{openproblem}
\textbf{Open Problem E.} Establish a closed-form connection between
$\PiP(r; k, \ell)$ and the Bateman--Horn local factor $\omega_f(\ell)$ for the
polynomials $f_{k, \eps}(m) = km(m+1) + \eps$. The numerical values of
$\PiP$ should be expressible as algebraic functions of $\omega_f(\ell; \eps)$.
\end{openproblem}

\begin{openproblem}
\textbf{Open Problem F.} Extend Theorems~\ref{thm:main} and~\ref{thm:degenerate}
to prime powers $\ell^j$. At $N = 10^6$, $k = 3$: residual biases appear at
modulus $27$ ($\chi^2 = 60.7$, $p < 10^{-4}$) and $81$ ($\chi^2 = 276$,
$p \ll 10^{-30}$). Theorems~\ref{thm:main} and~\ref{thm:degenerate} as stated
require $\ell$ prime; the prime-power extension is non-trivial.
\end{openproblem}

\section{Computational Methodology}
\label{sec:methodology}

All computations were performed on a standard laptop using Python 3.12 with
NumPy~\cite{numpy} and SciPy~\cite{scipy}. We document the algorithms here for
reproducibility; raw scripts and data will be deposited at Zenodo upon
publication.

\subsection{Database generation}

For each $k \in \{3, 5, 7, 9, 11, 13, 15, 17, 19, 21\}$, we generated the first
$N = 10^6$ triples $(m, \eps, q)$ representing the parametric family. The
algorithm proceeds in two stages.

\paragraph{Stage 1: Segmented Eratosthenes sieve.}
A boolean array $\sigma[0..P_{\max}]$ is initialised to \texttt{True}, then
sieved by setting $\sigma[i \cdot j] = \texttt{False}$ for each prime
$i \le \sqrt{P_{\max}}$ and $j \ge i$. For $P_{\max} = 1.8 \times 10^7$, this
generates all primes up to $P_{\max}$ in $\approx 0.3$ s.

\paragraph{Stage 2: Computation of $(m, \eps, q)$ from $p$.}
For each prime $p \ge 2k+1$ and each admissible $\eps \in \Eclass_k$, we test
whether $(p - \eps) / (2k)$ is an integer $T$. If so, we compute $m_{\mathrm{real}}
= (-1 + \sqrt{1 + 8T})/2$, and check the candidates $m \in \{\lfloor m_{\mathrm{real}}
\rfloor - 1, \lfloor m_{\mathrm{real}} \rfloor, \lfloor m_{\mathrm{real}} \rfloor + 1\}$
for the one giving smallest $|q| = |T - m(m+1)/2|$. The triple with smallest
$|q|$ across all $\eps$ is retained.

\paragraph{Verification.}
The reconstruction $p = km(m+1) + \eps + 2kq$ is verified for every triple.
Total generation time per database: $1.8$--$3.1$ s. Total disk: $\approx 5$ MB
per $k$ (compressed \texttt{.npz}).

\subsection{Statistical tests}

\paragraph{Pearson chi-square.}
For each $(k, \ell)$, we compute the histogram $\mathrm{obs}[r]$ of $q \bmod
\ell$ for $r \in \{0, \ldots, \ell-1\}$ and the test statistic $\chi^2 = \sum_r
(\mathrm{obs}[r] - N/\ell)^2 / (N/\ell)$, with $\ell - 1$ degrees of freedom
under $H_0$ (uniform distribution).

\paragraph{Permission frequency $\PiP$.}
For $(k, \ell)$ with $\gcd(2k, \ell) = 1$, we compute $(2k)^{-1} \bmod \ell$
via Python's \texttt{pow(2*k, -1, ell)} (modular inverse), then enumerate all
$(m_0, \eps) \in \Z/\ell\Z \times \Eclass_k$ to count $\#\{(m_0, \eps) :
q^*(k, m_0, \eps; \ell) = r\}$ for each $r$. Total time: $< 0.01$ s per case.

\paragraph{Effective constant $C_{\mathrm{eff}}$.}
For each subsample size $N \in \{10^2, 10^3, \ldots, 10^6\}$, the effective
constant is computed as $C_{\mathrm{eff}}(N) = \mathrm{err}_{\max}(N) \cdot \sqrt{N}$.

\subsection{Autocorrelation analysis}

\paragraph{ACF computation.}
For modulus $\ell$ and lag $j$, $\mathrm{ACF}_j = \mathrm{Corr}(s_n, s_{n+j})$
where $s_n = q_n \bmod \ell$, computed via NumPy's vectorised
\texttt{np.corrcoef(seq[:-j], seq[j:])}.

\paragraph{Wald--Wolfowitz runs test~\cite{wald}.}
We count the number of runs $R = 1 + \#\{n : s_n \neq s_{n+1}\}$ and compare
to its expected value $\mu_R = 2n_0 n_1/(n_0 + n_1) + 1$ under independence,
with variance $\sigma_R^2 = 2n_0 n_1 (2n_0 n_1 - n_0 - n_1) / ((n_0+n_1)^2
(n_0 + n_1 - 1))$, where $n_0, n_1$ are the counts of $0$s and $1$s in $s_n$.
The standardised statistic $z = (R - \mu_R)/\sigma_R$ is asymptotically
$\mathcal{N}(0, 1)$ under~$H_0$.

\paragraph{Damped cosine fit.}
Non-linear least squares fit (SciPy \texttt{optimize.curve\_fit}) of the form
$\alpha(\ell) = A \cdot e^{-\gamma(\ell-1)} \cos(\omega(\ell-1) + \varphi)$
with initial guess $(A, \omega, \varphi, \gamma) = (0.024, 1.5, 0, 0.5)$.

\subsection{Linear regressions}

For Conjecture~\ref{conj:acf}, the linear regression $\mathrm{ACF}_j(k) \approx
\alpha(j) k + \beta(j)$ is computed via SciPy \texttt{stats.linregress}, which
returns slope, intercept, $R^2$, and $p$-value (two-sided test against zero
slope).

\subsection{Computational cost summary}

\begin{table}[h!]
\centering
\caption{Summary of computational costs.}
\begin{tabular}{l c}
\toprule
Operation & Time \\
\midrule
Generate $10^6$ triples for one $k$ & $1.8$--$3.1$ s \\
$\chi^2$ tests for one $(k, \ell)$ pair & $< 0.01$ s \\
Compute $\PiP$ for one $(k, \ell)$ & $< 0.01$ s \\
ACF lag-1 to lag-10 for one $k$ & $\approx 1$ s \\
Damped cosine fit (10 data points) & $< 0.01$ s \\
Bivariate regression (Section~\ref{sec:anti-persistence}) & $< 0.1$ s \\
\midrule
\textbf{Full Section~\ref{sec:validation-k3}+\ref{sec:universality} analysis} & $\approx 60$ s \\
\textbf{Full Section~\ref{sec:anti-persistence} (ACF) analysis} & $\approx 5$ s \\
\bottomrule
\end{tabular}
\end{table}

\subsection{Reproducibility}

The complete pipeline (sieve, $(m, \eps, q)$ extraction, statistical tests,
ACF analysis, regression fits, figure generation) runs in under two minutes
on a standard laptop with $\le 2$ GB RAM. The Python source code and
compressed \texttt{.npz} databases (one per $k$) will be deposited at Zenodo.
Each database contains four arrays: $p$, $m$, $\eps$, $q$, each of length
$10^6$.

\section{Summary Epistemic Table}
\label{sec:epistemic}

\begin{table}[h!]
\centering
\caption{Complete epistemic status of all results.}
\begin{tabular}{l l p{5.5cm}}
\toprule
Result & Status & Note \\
\midrule
Proposition~\ref{prop:forbidden} (forbidden residue) & Rigorous & 3-line proof \\
Theorem~\ref{thm:degenerate} (degenerate case) & Rigorous & Under H1$'$ (BV at fixed modulus) \\
Theorem~\ref{thm:main} (non-degenerate case) & Heuristic & Under H1$'$+H3$'$ \\
Corollary~\ref{cor:constant} (explicit $C(k, \ell)$) & Rigorous (under~\ref{thm:main}) & Triangle bound \\
Proposition~\ref{prop:h3-implies-h3prime} (H3 $\Rightarrow$ H3$'$) & Heuristic & 5-step proof, Section~\ref{sec:connection} \\
Conjecture~\ref{conj:tight} (optimal $C \approx 0.5$) & Open & Validated empirically \\
Conjecture~\ref{conj:acf} (ACF linear law + cosine) & Open & New empirical phenomenon, $R^2 = 0.999$ \\
Observation~\ref{disc:cosine} (damped cosine signature) & Discovery & 10 data points, $R^2 = 0.9988$ \\
H1$'$ & Heuristic & Bombieri--Vinogradov consequence \\
H3$'$ & Conjectural & Analog of H3 of Paper~II \\
H3 (Paper~II) $\Rightarrow$ H3$'$ & Heuristic (proved) & 5-step proof in~\S\ref{sec:connection} \\
\bottomrule
\end{tabular}
\end{table}

\section{Conclusions}

This paper completes the trilogy of Papers I--III on the parametric prime
family $p = km(m+1) + \eps + 2kq$, providing the modular layer of description
that complements the marginal (Paper~I, $\E[|\qmin|]$) and conditional
(Paper~II, $\E[|q|\mid m]$) layers.

\paragraph{Three principal results.}
\begin{enumerate}[leftmargin=*]
\item An \textbf{unconditional theorem} (Theorem~\ref{thm:degenerate}) for
the degenerate case $\ell \mid 2k$: $F_N(r; k, \ell) = 1/\ell + O(N^{-1/2})$.
Validated across all 10 cases tested with $\chi^2 \le 3.1$.
\item A \textbf{conditional theorem} (Theorem~\ref{thm:main}) for the
non-degenerate case, with explicit forbidden residue
(Proposition~\ref{prop:forbidden}) and explicit effective constant
(Corollary~\ref{cor:constant}). The empirical effective constant is
$\bar C \approx 0.5$, $5\times$ smaller than the triangle-inequality bound
(Conjecture~\ref{conj:tight}).
\item A \textbf{new empirical phenomenon}
(Discovery~\ref{disc:acf}--\ref{disc:cosine}): the lag-1 autocorrelation of
$(q_n \bmod 2)$ obeys a clean linear law $\mathrm{ACF}_1 \approx 0.024\,k - 0.20$
($R^2 = 0.97$) with sign-reversal threshold at $k_0 \approx 8.4$. The full
spectral signature $\{\alpha(\ell)\}_{\ell=1}^{10}$ of the slopes per lag is
fitted by a damped cosine (period $\approx 5.89$, damping $\approx 0.39$)
to within $R^2 = 0.999$.
\end{enumerate}

\paragraph{Strategic connection to Paper II.} Hypothesis H3$'$ is a structural
analogue of Paper~II's H3. Theorem~\ref{thm:main} provides an indirect
quantitative test of H3: the empirical agreement ($r > 0.99$ for small $\ell$)
is consistent with H3 holding in the regime tested. We are careful not to
overstate this: numerical evidence at $N = 10^6$, however clean, cannot settle
a Bateman--Horn-type conjecture, as the historical record reminds us
(\textit{cf}.\ the Hardy--Littlewood incompatibilities; Skewes-type
sign-changes invisible up to $10^{316}$). The role of
Proposition~\ref{prop:h3-implies-h3prime} is therefore precisely the role we
claim for it: to establish that an unconditional proof of H3 (Paper~II) would
imply our predictions, making our predictions falsifiable consequences of a
deeper unsolved problem.

\paragraph{Honest assessment of difficulty.} An unconditional proof of
Theorem~\ref{thm:main} is currently out of reach for three independent
reasons (Section~\ref{sec:connection}): (i) Bateman--Horn itself is open
beyond the linear case; (ii) the parity barrier blocks all known sieve
approaches to primes represented by polynomials of degree $\ge 2$;
(iii) uniform PNT in arithmetic progressions of moduli $\sim N^{1/2}$ is
unavailable. Theorem~\ref{thm:degenerate} is the only genuinely unconditional
content of the paper; it is rigorous because it reduces to the linear
(Dirichlet) case. We also caution that the empirical signature of
Section~\ref{sec:anti-persistence}---the damped cosine fit at $R^2 = 0.999$---
may itself be the trace of a non-trivial second-order correlation that the
Cramér--Granville hypothesis (H2) ignores; this is precisely the content of
Open Problem~D.

\paragraph{New research directions.} Beyond the open problems explicitly
listed (Section~\ref{sec:openpb}), three directions seem natural for future
work:
\begin{enumerate}[leftmargin=*]
\item \textbf{The damped-cosine origin}. The fact that the sequence
$\{\alpha(\ell)\}$ admits a damped-cosine fit to $R^2 = 0.999$ on 10 data
points strongly suggests an underlying analytic mechanism. Identifying it---
for instance via a transfer-operator analysis of the joint distribution of
$(\qmin)_n$ and $(\qmin)_{n+1}$---would constitute a substantial advance.
\item \textbf{Modulus dependence of the spectral parameters}. Whether the
period $T \approx 5.89$ and damping rate $\gamma \approx 0.39$ persist for
$\ell > 2$, or vary with $\ell$ in a structured way, is an immediate empirical
question that may shed light on the mechanism.
\item \textbf{Prime power moduli}. Open Problem~F notes residual biases at
$\ell^j$ for $\ell^j \in \{27, 81\}$. The full extension of
Theorem~\ref{thm:main} to prime powers requires reformulating
Proposition~\ref{prop:forbidden} for non-prime $\ell$, and is non-trivial.
\end{enumerate}

\paragraph{Negative result and methodological caveat.} The pilot study at
$N = 10^4$ for $k \ge 7$ was contaminated by a window-truncation artefact;
validation at $N = 10^6$ was required to distinguish genuine $p$-adic biases
from geometric artefacts. This methodological lesson is consistent with the
spurious-regression caveat established in Paper~I, and reinforces the general
principle that scale-dependent effects in number-theoretic data require
large-$N$ validation.

\paragraph{Final remark.} Our results illustrate the fertile interplay
between explicit number-theoretic computation, statistical analysis, and
analytic conjectures. Each of the three layers (marginal, conditional,
modular) of the parametric prime family is now characterised by a precise
quantitative law, validated to high precision, with carefully tracked
hypotheses. The remaining open problems (A--F) point to specific places
where progress on classical questions (parity barrier, large-modulus
distribution) would have direct, falsifiable consequences for our predictions.

\section*{Acknowledgements}
The author thanks the Python open-source community (NumPy, SciPy) for making
large-scale computational number theory accessible. All computations were
performed on a standard laptop; all $10^6$-prime databases were generated
in under 3 seconds each.

\section*{Data Availability}

All numerical data and Python scripts will be deposited at Zenodo upon
acceptance. The current working version is available upon request.

\begin{thebibliography}{40}

% === Vos propres papiers (auto-citation) ===
\bibitem{paper1} H.~Bakkaoui, \textit{A Parametric Family of Primes via Quadratic
Forms: Heuristic Distribution of the Minimal Offset, Large-Scale Numerical
Validation, and Statistical Control of Spurious Spectral Correlations}, Working
Paper~I, April 2026.

\bibitem{paper2} H.~Bakkaoui, \textit{Parametric Prime Families via Quadratic
Forms: Layer Invariants, Conditional Proofs, and Connections to the Riemann
Zeta Function}, Working Paper~II (revised), April 2026.

% === Théorèmes classiques de la théorie analytique des nombres ===
\bibitem{dirichlet} P.~G.~L.~Dirichlet, \textit{Beweis des Satzes, dass jede
unbegrenzte arithmetische Progression unendlich viele Primzahlen enth\"{a}lt},
Abh.\ K\"{o}n.\ Preuss.\ Akad.\ Wiss.\ (1837), 45--81.

\bibitem{batemanhorn} P.~T.~Bateman and R.~A.~Horn, \textit{A heuristic
asymptotic formula concerning the distribution of prime numbers}, Math.\ Comp.\
\textbf{16} (1962), 363--367.

\bibitem{aletheia} S.~L.~Aletheia-Zomlefer, L.~Fukshansky, and S.~R.~Garcia,
\textit{The Bateman-Horn conjecture: heuristics, history, and applications},
Expositiones Math.\ \textbf{38} (2020), 430--479.

\bibitem{garcia} S.~R.~Garcia, \textit{What is the Bateman--Horn conjecture?},
Notices Amer.\ Math.\ Soc.\ \textbf{71} (2024), 1378--1381.

\bibitem{bombieri} E.~Bombieri, \textit{On the large sieve}, Mathematika
\textbf{12} (1965), 201--225.

\bibitem{vinogradov-bv} A.~I.~Vinogradov, \textit{The density hypothesis for
Dirichlet $L$-series}, Izv.\ Akad.\ Nauk SSSR Ser.\ Mat.\ \textbf{29} (1965),
903--934.

\bibitem{cramer} H.~Cram\'{e}r, \textit{On the order of magnitude of the
difference between consecutive prime numbers}, Acta Arith.\ \textbf{2} (1936),
23--46.

\bibitem{granville} A.~Granville, \textit{Harald Cram\'{e}r and the distribution
of prime numbers}, Scand.\ Actuar.\ J.\ (1995), 12--28.

\bibitem{tao} T.~Tao, \textit{Cram\'{e}r's random model and related models for
the primes}, Blog post, January 2015.

\bibitem{hardylittlewood} G.~H.~Hardy and J.~E.~Littlewood, \textit{Some
problems of Partitio Numerorum III: on the expression of a number as a sum of
primes}, Acta Math.\ \textbf{44} (1923), 1--70.

\bibitem{rosser} J.~B.~Rosser and L.~Schoenfeld, \textit{Approximate formulas
for some functions of prime numbers}, Illinois J.\ Math.\ \textbf{6} (1962),
64--94.

\bibitem{cipolla} M.~Cipolla, \textit{La determinazione asintotica dell'$n$-esimo
numero primo}, Rend.\ Accad.\ Sci.\ Fis.\ Mat.\ Napoli \textbf{8} (1902), 132--166.

% === Crible et bornes effectives ===
\bibitem{iwaniec} H.~Iwaniec and E.~Kowalski, \textit{Analytic Number Theory},
AMS Colloquium Publications, Vol.~53, 2004.

\bibitem{davenport} H.~Davenport, \textit{Multiplicative Number Theory}, 3rd ed.\
(revised by H.~L.~Montgomery), Springer, 2000.

\bibitem{heathbrown} D.~R.~Heath-Brown, \textit{Zero-free regions for Dirichlet
$L$-functions, and the least prime in an arithmetic progression}, Proc.\ London
Math.\ Soc.\ (3) \textbf{64} (1992), 265--338.

\bibitem{linnik} U.~V.~Linnik, \textit{On the least prime in an arithmetic
progression. I. The basic theorem}, Mat.\ Sbornik N.S.\ \textbf{15(57)} (1944),
139--178.

\bibitem{harman} G.~Harman, \textit{Primes in short intervals}, Math.\ Z.\
\textbf{180} (1982), 335--348.

\bibitem{bhp} R.~C.~Baker, G.~Harman, and J.~Pintz, \textit{The difference
between consecutive primes, II}, Proc.\ London Math.\ Soc.\ \textbf{83} (2001),
532--562.

\bibitem{maynard} J.~Maynard, \textit{Small gaps between primes}, Ann.\ of
Math.\ \textbf{181} (2015), 383--413.

\bibitem{selberg} A.~Selberg, \textit{On an elementary method in the theory of
primes}, Norske Vid.\ Selsk.\ Forh., Trondhjem \textbf{19} (1947), 64--67.

\bibitem{friedlander} J.~Friedlander and H.~Iwaniec, \textit{Opera de Cribro},
AMS Colloquium Publications, Vol.~57, 2010.

% === Fonction zeta et hypothèse de Riemann ===
\bibitem{riemann} B.~Riemann, \textit{\"Uber die Anzahl der Primzahlen unter
einer gegebenen Gr\"osse}, Monatsberichte der Berliner Akademie (1859), 671--680.

\bibitem{montgomery} H.~L.~Montgomery, \textit{The pair correlation of zeros of
the zeta function}, In: Analytic Number Theory, Proc.\ Symp.\ Pure Math.\
\textbf{24} (1973), 181--193.

\bibitem{odlyzko} A.~M.~Odlyzko, \textit{On the distribution of spacings between
zeros of the zeta function}, Math.\ Comp.\ \textbf{48} (1987), 273--308.

\bibitem{titchmarsh} E.~C.~Titchmarsh, \textit{The Theory of the Riemann
Zeta-Function}, 2nd ed., Oxford University Press, 1986.

% === Biais de Chebyshev et distribution dans les progressions arithmétiques ===
\bibitem{rubinstein} M.~Rubinstein and P.~Sarnak, \textit{Chebyshev's bias},
Experiment.\ Math.\ \textbf{3} (1994), 173--197.

\bibitem{fiorillimartin} D.~Fiorilli and G.~Martin, \textit{Inequities in the
Shanks-R\'{e}nyi prime number race: an asymptotic formula for the densities},
J.\ Reine Angew.\ Math.\ \textbf{676} (2013), 121--212.

\bibitem{granville-martin} A.~Granville and G.~Martin, \textit{Prime number
races}, Amer.\ Math.\ Monthly \textbf{113} (2006), 1--33.

% === Tests numériques et primalité ===
\bibitem{rabin} M.~O.~Rabin, \textit{Probabilistic algorithm for testing
primality}, J.\ Number Theory \textbf{12} (1980), 128--138.

\bibitem{psw} C.~Pomerance, J.~L.~Selfridge, and S.~S.~Wagstaff Jr., \textit{The
pseudoprimes to $25 \cdot 10^9$}, Math.\ Comp.\ \textbf{35} (1980), 1003--1026.

\bibitem{crandall} R.~Crandall and C.~Pomerance, \textit{Prime Numbers: A
Computational Perspective}, 2nd ed., Springer, 2005.

\bibitem{atkin} A.~O.~L.~Atkin and F.~Morain, \textit{Elliptic curves and
primality proving}, Math.\ Comp.\ \textbf{61} (1993), 29--68.

% === Statistiques (tests permutation, séries temporelles) ===
\bibitem{wald} A.~Wald and J.~Wolfowitz, \textit{On a test whether two samples
are from the same population}, Ann.\ Math.\ Statist.\ \textbf{11} (1940),
147--162.

\bibitem{theiler} J.~Theiler, S.~Eubank, A.~Longtin, B.~Galdrikian, and J.~D.~Farmer,
\textit{Testing for nonlinearity in time series: the method of surrogate data},
Physica D \textbf{58} (1992), 77--94.

\bibitem{granger} C.~W.~J.~Granger and P.~Newbold, \textit{Spurious regressions
in econometrics}, J.\ Econometrics \textbf{2} (1974), 111--120.

\bibitem{hurst} H.~E.~Hurst, \textit{Long-term storage capacity of reservoirs},
Trans.\ Amer.\ Soc.\ Civil Eng.\ \textbf{116} (1951), 770--799.

% === Familles de premiers et applications ===
\bibitem{cox} D.~A.~Cox, \textit{Primes of the Form $x^2 + ny^2$: Fermat, Class
Field Theory, and Complex Multiplication}, Wiley-Interscience, 1989.

\bibitem{jacobson} M.~J.~Jacobson Jr.\ and H.~C.~Williams, \textit{New quadratic
polynomials with high densities of prime values}, Math.\ Comp.\ \textbf{72}
(2003), 499--519.

\bibitem{li-bh} W.~Li, \textit{A note on the Bateman--Horn conjecture}, J.\
Number Theory \textbf{208} (2020), 390--399.

\bibitem{weyl} H.~Weyl, \textit{\"Uber die Gleichverteilung von Zahlen mod.\
Eins}, Math.\ Ann.\ \textbf{77} (1916), 313--352.

\bibitem{vinogradov-sums} I.~M.~Vinogradov, \textit{Some theorems concerning
the theory of primes}, Mat.\ Sbornik N.S.\ \textbf{2(44)} (1937), 179--194.

% === Outils computationnels ===
\bibitem{numpy} C.~R.~Harris et al., \textit{Array programming with NumPy},
Nature \textbf{585} (2020), 357--362.

\bibitem{scipy} P.~Virtanen et al., \textit{SciPy 1.0: fundamental algorithms
for scientific computing in Python}, Nature Methods \textbf{17} (2020), 261--272.

\end{thebibliography}

\end{document}
